\section*{3. 对 Lefschetz 线束的应用}

\subsubsection*{3.1.}
从现在起, 固定
\begin{equation}\tag{3.1.0}
k,\quad\text{为代数闭域},
\end{equation}
以及
\begin{equation}\tag{3.1.1}
\begin{aligned}
X\subset\mathbb P^r,\quad
&\text{射影、光滑且连通的 }k\text{-概形}\\
&\text{偶数维数为 }n\text{, 满足 (LV) (XVIII~5.2.2).}
\end{aligned}
\end{equation}
取
\begin{equation}\tag{3.1.2}
D=\{H_t\}_{t\in\mathbb P^1},
\quad\text{次数为 $d$ 的超曲面的 Lefschetz 线束},
\end{equation}
其泛纤维为
\begin{equation}\tag{3.1.3}
X_\eta,\qquad \eta=\text{为下列空间的泛点: }\mathbb P^1.
\end{equation}
置 (参见~XVIII~5.4.4)
\begin{equation}\tag{3.1.4}
m:X_{\bar\eta}\hookrightarrow X\times_k k(\bar\eta),
\quad\text{为包含映射},
\end{equation}
以及
\begin{equation}\tag{3.1.5}
E^{n-1}(X_{\bar\eta},\mathbb Q_\ell(i))
=\bigl(m^*H^{n-1}(X,\mathbb Q_\ell(i))\bigr)^\perp
\subset H^{n-1}(X_{\bar\eta},\mathbb Q_\ell(i)).
\end{equation}
在 XVIII~6.7.2 中已见, \underline{Picard--Lefschetz 公式}蕴含\footnote{印本在 (3.1.6) 中写作 $E^{2n-1}$. (3.1.5) 中定义且命题~3.2 的证明中始终使用的空间为 $E^{n-1}$.}
\begin{equation}\tag{3.1.6}
\begin{gathered}
\text{若 }n\text{ 为偶数, 则群 }
\operatorname{Gal}(k(\bar\eta)/k(\eta))\\
\text{的任意有限指数子群均不可约地作用于 }E^{n-1}(X_{\bar\eta},\mathbb Q_\ell(i)).
\end{gathered}
\end{equation}

\subsubsection*{命题 3.2.}
设 $D=\{H_t\}$ 为 Lefschetz 线束, 并假定 $n=2m$ 为偶数. 则或者有
\begin{equation}\tag{3.2.1}
H^{n-1}(X_{\bar\eta},\mathbb Q_\ell(i))
\subset m^*H^{n-1}(X,\mathbb Q_\ell(i))
+N^{m-1}H^{n-1}(X_{\bar\eta},\mathbb Q_\ell(i)),
\end{equation}
或者有:

\subsubsection*{(3.2.2)}
对于每条在 $k(\bar\eta)$ 上固有且光滑的曲线 $C$ 和 $C\times X_{\bar\eta}$ 上每个余维为 $m$ 的代数圈 $z$, 所诱导的同态 (1.2.3)
\begin{equation}\tag{3.2.3}
\mathrm{cl}(z):H^1(C,\mathbb Q_\ell(1+i-m))
\longrightarrow H^{2m-1}(X_{\bar\eta},\mathbb Q_\ell(i))
\end{equation}
的像均位于 $m^*H^{2m-1}(X,\mathbb Q_\ell(i))$ 中.

\noindent\underline{证明.}
假定某个这样的同态 (3.2.3) 的像不包含于 $m^*H^{n-1}(X,\mathbb Q_\ell(i))$ 中, 或等价地, 它在 $E^{n-1}(X_{\bar\eta},\mathbb Q_\ell(i))$ 中的像非零. 显然, 此像在 $\operatorname{Gal}(k(\bar\eta)/k(\eta))$ 的某个有限指数子群的作用下稳定. 因而 $E^{n-1}(X_{\bar\eta},\mathbb Q_\ell(i))$ 的不可约性 (3.1.6) 表明
\begin{equation*}
H^{n-1}(X_{\bar\eta},\mathbb Q_\ell(i))
=m^*H^{n-1}(X,\mathbb Q_\ell(i))+\operatorname{Im}(3.2.3).
\end{equation*}
根据定义,
\begin{equation*}
\operatorname{Im}(3.2.3)
=\mathrm{cl}(z)\bigl(H^1(C,\mathbb Q_\ell(1+i-m))\bigr)
\subset N^{m-1}H^{n-1}(X_{\bar\eta},\mathbb Q_\ell(i)).
\end{equation*}
所以,
\begin{equation*}
H^{n-1}(X_{\bar\eta},\mathbb Q_\ell(i))
=m^*H^{n-1}(X,\mathbb Q_\ell(i))
+N^{m-1}H^{n-1}(X_{\bar\eta},\mathbb Q_\ell(i)),
\end{equation*}
这正是第一种情形.
